% =============================================================================
% §2 Background  — Target: 0.8--1.0 page. No new claims.
% 4.5/5 bar: enough for a non-ANNS ASPLOS reviewer to follow Design/Eval;
% every compared system is named with its medium and interface.
% 填写: 补 CXL-SSD 设备框图（Fig.~2）上的真实端口/窗口数字。
% =============================================================================
\section{Background and Deployment Points}
\label{sec:background}

\subsection{Graph ANNS and the block-SSD path}
\label{sec:bg-anns}

A graph ANNS index (Vamana~\cite{diskann}, HNSW~\cite{hnsw}) stores, for each
vector, a small neighbor list. Search maintains a candidate set of width
$L$ (beam / \texttt{efSearch}), repeatedly expanding the closest unread node,
and returns the top-$k$. Quality is monotone in how many nodes are scored;
\emph{stall} is monotone in how many of those scores miss the fast tier.

DiskANN~\cite{diskann,freshdiskann} is the reference block-SSD design. Full-
precision vectors and adjacency live in 4\,KB-aligned sectors on NVMe. A
compact PQ codebook in DRAM scores neighbors without a NAND read; an
\texttt{aio} beam issues several sector reads at once; a static node cache
pins entry vertices. PipeANN~\cite{pipeann} keeps that layout and overlaps
compute with I/O inside a single query, and can interleave queries to raise
device queue depth. The first-class objects of this path are \emph{the I/O
queue and the sector}. They are the wrong first-class objects once the
device is no longer a block target but a faulting address space
(\S\ref{sec:motivation}).

\subsection{CXL-DRAM ANNS}
\label{sec:bg-cxl-dram}

CXL.mem lets a host map device HDM into its physical address space~\cite{cxl30}.
CXL-ANNS~\cite{cxlanns,cxlanns-tocs} places the essential graph and embeddings
in a Type-3 CXL-DRAM pool and hides $\sim$100\,ns far-memory delay with a
relationship-aware cache and DSA-assisted prefetch. Cosmos~\cite{cosmos}
further offloads traversal and distance to device-side compute. Second-tier
and residual-quantization designs~\cite{secondtier,fatrq,hm-ann} also assume
a DRAM- or NVM-priced far tier. \sys{} shares the \emph{interface}
(load/store) with this line and rejects its \emph{capacity assumption}
(the corpus fits in CXL-DRAM).

\subsection{Flash-backed CXL memory}
\label{sec:bg-cxl-ssd}

A CXL-SSD pairs NAND with a modest DRAM cache and presents a byte-addressable
window~\cite{skybyte,fromblocktobyte,cmm-h}. A hit is a CXL.mem load; a miss
is a NAND fill into that window, then a retry. The programming model looks
like memory; the miss tax looks like flash. Hardware and generic OS work on
this tier does not manage graph-ANNS residency. CXL-AnySSD~\cite{cxl-anyssd}
accelerates vector-DB traffic by removing swap, not by controlling promote
precision or search stalls. P-HNSW~\cite{phnsw} studies crash-consistent
graphs and mentions CXL-SSD; its experiments use DRAM-simulated PM.

Three conditions must not be mixed when we talk about sharing
(Figure~\ref{fig:bg-cxl}):

\begin{enumerate}
  \item \textbf{Direct-attach CXL-SSD.} One host, flash price, memory
    semantics. Sufficient for the single-host claims in this paper.
  \item \textbf{CXL 2.0 pooling / MLD.} Several hosts may time-share a
    device; a given HDM region is owned by one host at a time.
  \item \textbf{CXL 3.x Shared FAM / GFAM.} Several hosts may map the same
    region at once. This is the standard premise for multi-host one-copy
    serving, and additionally requires a fabric manager and a read-only
    publish protocol. We do not evaluate it.
\end{enumerate}

\begin{figure}[t]
  \centering
  \phbox{0.98\columnwidth}{3.2cm}{%
    Fig.~2 placeholder --- CXL-SSD stack:\\
    host CPU $\rightarrow$ CXL link $\rightarrow$ endpoint
    (DRAM window $+$ FTL $+$ NAND).\\
    Annotate: host-side link is per-host; device-side bandwidth is shared.}
  \caption{Flash-backed CXL memory. Host-side CXL links are independent;
    device-side port, controller, DRAM cache, FTL, and NAND are shared.
    One-copy reduces replicated capacity; it does not replicate SSD IOPS.
    \textbf{Fill-in:} use the real endpoint you measure; mark the software
    proxy path if that is the prototype.}
  \Description{Block diagram of a host attached to a CXL-SSD endpoint with
    a small DRAM window in front of NAND.}
  \label{fig:bg-cxl}
\end{figure}

\paragraph{Media-cost sketch (qualitative).}
Let $c_{\mathrm{DRAM}}$ and $c_{\mathrm{flash}}$ be unit media prices and
$n_{\mathrm{HA}}$ the number of high-availability copies ($n_{\mathrm{HA}}\ll R$
in the common case). Then
\[
  \mathrm{MediaCost}_{\text{CXL-DRAM}} \sim D\cdot c_{\mathrm{DRAM}},
\]
\[
  \mathrm{MediaCost}_{\text{CXL-SSD}} \sim n_{\mathrm{HA}} D\cdot c_{\mathrm{flash}}
    + H\cdot c_{\mathrm{cache}} + C_{\mathrm{fabric}}.
\]
We use this only to locate the third deployment point
(Figure~\ref{fig:eval-scale}). We do not claim a lower TCO.
